\begin{center}\large\bfseries NACHTRAG.\end{center}
Nach dem Druck wurde in der Vergleichstabelle an zwei Orten ein kleiner Fehler bemerkt. In Callao geht er auf eine falsche Längenangabe in der angeführten Schrift zurück, auf St. Helena auf einen Rechenfehler. Ich nutze diese Gelegenheit, um hier zusammen mit den Ergebnissen einer verbesserten Rechnung auch die Theorie mit den Beobachtungen an acht weiteren Orten zu vergleichen, von denen ich inzwischen erfahren habe. [Die Berichtigungen sind im Wiederabdruck berücksichtigt. Zur besseren Übersicht wurde der Vergleich der Beobachtungen an diesen acht Orten mit denen an den ursprünglich 91 Orten bereits oben zusammengestellt.]
\SourceLeaf{180}
Die Beobachtungen in Stockholm stammen von \textsc{Rudberg}, Intensität und Inklination aus dem Jahr 1832, die Deklination aus dem Jahr 1833 (\textsc{Poggendorff}s Annalen, Band 37). Die Beobachtungen in Brüssel stammen aus dem Jahr 1832. Die Deklination und Inklination wurden von \textsc{Quetelet} bestimmt (Bulletins de l’Académie de Bruxelles, Band VI), die Intensität von \textsc{Rudberg} (die oben [S. 154] angeführte Schrift von \textsc{Sabine}). Der freundlichen Mitteilung \textsc{Sabine}s verdanke ich die Bestimmungen für die übrigen neuen Orte sowie die Bestimmung der Intensität und eine neuere Beobachtung der Inklination in Callao. Die Beobachtungen in Lerwick und Valentia wurden 1838 von Kapitän \textsc{James Ross} angestellt, die in Port Etches, Panama und Oahu 1837 von Kapitän \textsc{Belcher}. die in Callao 1838 von demselben. In Montreal wurden schließlich die Inklination und Intensität 1838 von Major \textsc{Estcourt} beobachtet. Die Deklination stammt dagegen aus dem Jahr 1834, und der Beobachter wird nicht genannt.

Zur Bildtafel, die die in Artikel 12 entwickelten Untersuchungen veranschaulicht, ist hier noch zu bemerken, dass der geschickte Lithograf Herr \textsc{Rittmüller} versucht hat, zugleich die unterschiedliche Intensität darzustellen. Er tat dies auf zwei Arten, durch die unterschiedliche Stärke der Linien und die ungleiche Schattierung der Zwischenräume.

Weil sich der Abschluss des Drucks dieses Bandes verzögert hat, konnten ihm außer der Karte für die Werte von $V$ [siehe Artikel 30] noch zwei weitere Karten beigefügt werden. Die erste, die die nach den Elementen oder aus den Tabellen \textit{berechneten} Werte der Deklination darstellt, verdanken die Leser meinem verehrten Freund und Mitherausgeber der \textit{Resultate}. Um die verwickelte Anordnung des Systems der Linien gleicher Deklination deutlich überblicken zu können, wurden die Punkte, an denen die Deklination einen Maximalwert hat, ebenso wie die Punkte, an denen sich zwei Linien gleicher Deklination kreuzen (oder eine sich selbst kreuzt), mit besonderer Sorgfalt berechnet. Punkte der ersten Art gibt es zwei, Punkte der zweiten Art vier. Diesen Punkten ist gemeinsam, dass an ihnen das erste Differential der Deklination in jeder Richtung verschwindet. Es ist übrigens überflüssig zu bemerken, dass in Gegenden, in denen sich die Deklination in alle Richtungen langsam ändert, wie im südlichen und südöstlichen Asien, schon geringe Änderungen der Deklinationswerte sehr große Änderungen in der Anordnung der Linien bewirken können.
